\documentclass[11pt]{scrartcl}

%%%%%%%%%%%%%% LATEX SAMPLE FILE %%%%%%%%%%%%%%%%
% A line which starts with a % sign
% is called a COMMENT. It is IGNORED
% by the LaTeX processor.

% Include math
\usepackage{amsmath,amsthm,amssymb,graphicx,enumitem}
% Include links
\usepackage[hidelinks]{hyperref}

%%%%%%%%%%%%%  THEOREMS  %%%%%%%%%%%%%%%%
\theoremstyle{plain}
\newtheorem{theorem}{Theorem}
\newtheorem{lemma}[theorem]{Lemma}
\theoremstyle{definition}
\newtheorem*{definition}{Definition}

\theoremstyle{remark}
\newtheorem{remark}{Remark}

% Problem command
\newcounter{problem}
\newcommand{\problem}[2]{%
    \refstepcounter{problem}%
    \noindent\textbf{\href{#1}{Problem~\theproblem.}}#2%
    \vspace{0.3cm}\par%
}

% Solution command
\newcounter{solution}
\newcommand{\solution}[1]{%
    \refstepcounter{solution}%
    \noindent\textbf{Solution.~}#1%
    \vspace{0.3cm}\par%
}

% Solution command that is used when there are more than one solutions.
\newcounter{numsolution}
\newcommand{\numsolution}[2]{%
    \refstepcounter{numsolution}%
    \noindent\textbf{Solution~#1.~}#2%
    \vspace{0.3cm}\par%
}

\renewcommand{\mod}{\text{mod }}
\newcommand\cycle[2][\,]{%
\readlist\thecycle{#2}%
(\foreachitem\i\in	hecycle{\ifnum\icnt=1\else#1\fi\i})%
}

%%%%%%%%%%%%%%  PAGE SETUP %%%%%%%%%%%%%%%%
\usepackage[top=1in,left=1in,right=1in,bottom=1.5in]{geometry}
\usepackage{fancyhdr}
\pagestyle{fancy}
\lhead{DAA Project}
\rhead{\thepage}
\cfoot{}

%%%%%%%%%%%%% SHORTCUTS %%%%%%%%%%%%%%%%%%%%
\newcommand{\half}{\frac{1}{2}}
\newcommand{\cbrt}[1]{\sqrt[3]{#1}}

\begin{document}
\title{Design and Analysis of Algorithms Project Report}
\author{Anshul Raj Singh, Gauransh Kapur}
\date{\today}
\maketitle

\tableofcontents

\section{Problem statement and test cases}
One day, Alex got tired of climbing in a gym and decided to take a very large group of climbing friends outside to climb. The climbing area where they went, had a huge wide boulder, not very tall, with various marked hand and foot holds. Alex quickly determined an ``allowed" set of moves that her group of friends can perform to get from one hold to another. The overall system of holds can be described by a rooted tree
$T$ with $n$ vertices, where each vertex corresponds to a hold and each edge corresponds to an allowed move between holds. The climbing paths converge as they go up the boulder, leading to a unique hold at the summit, represented by the root of $T$. Alex and her friends (who are all excellent climbers) decided to play a game, where as many climbers
as possible are simultaneously on the boulder and each climber needs to perform a sequence of exactly $k$ moves. Each climber can choose an arbitrary hold to start from, and all moves must move away from the ground. Thus, each climber traces out a path of $k$ edges in the tree $T$, all directed toward the root. However, no two climbers are allowed to
touch the same hold; the paths followed by different climbers cannot intersect at all. 

Describe and analyze an efficient dynamic programming algorithm to compute the maximum number of climbers that can play this game.




\section{Solution}
Let $r$ be the root of the tree $T$. We define two quantities for each vertex $v$ in $T$.\\
\begin{itemize}
    \item Let $G(v)$ be the maximum number of $k$-length paths that can be chosen in the subtree rooted at $v$ not containing $v$.
    \item Let $H(v, d)$ be the maximum number of $k$-length paths that can be chosen in the subtree rooted at $v$ subject to the restriction that there is a path (call it \textbf{distinguished}) at depth $d$ containing $v$ for $1\leq d\leq k$
\end{itemize}Then the required answer is \[\max\{G(r), H(r, k)\}.\]
So we need to compute $G(v)$ and $H(v, d)$ for all vertices $v$ in $T$ and for all $1\leq d\leq k$. We proceed via recursion.\\

\subsection{Recurrence for $G(v)$}
Let $u_1, \ldots, u_m$ be the children of $v$. For every $1\leq j\leq m$, the maximum number of $k$-length paths that end at $u_j$ is $\max\{G(u_j), H(u_j, k)\}$. Thus, \[G(v) = \sum_{j=1}^m \max\{G(u_j), H(u_j, k)\}.\]

\subsection{Recurrence for $H(v, d)$}
For $1\leq d < k$, let $u$ be that child of $v$ such that the distinguished path goes through $u$. Then, the maximum number of $k$-length paths that can be chosen in the subtree, keeping $u$ fixed, is \[H(u, d-1) + \sum_{\substack{w \text{ child of } v\\ w\neq u}}\max\{G(w), H(w, k)\}.\]Therefore, we have the following recurrence \[H(v, d) = \max_{u \text{ child of } v} \left(H(u, d-1) + \sum_{\substack{w \text{ child of } v\\ w\neq u}}\max\{G(w), H(w, k)\}\right)\] for $1\leq d < k$.

When $d = k$, then the distingueshed path finishes at $v$, so this path must be counted in the final sum. This gives \[H(v, k) = 1 + \max_{u \text{ child of } v}\left(H(u, d-1) + \sum_{\substack{w \text{ child of } v\\ w\neq u}}\max\{G(w), H(w, k)\}\right).\]So, our final answer is $$\max\{G(r), H(r, k)\}.$$

\section{Python Code}



\section{Analysis of Algorithm}

\section{Conclusion}

\end{document}